\begin{pastexam}[label={px:nsysu-114-dmalg-05}]{中山資工所資安組 114 學年度 第 6 題}
（計算）

{[10\%]} If a fair die is rolled for 15 times, what is the probability that the sum of the rolls
is 43?
\end{pastexam}

\begin{pastsoln}
答案：
\[
\frac{1}{6^{15}} \sum_{k=0}^{4} (-1)^k \binom{15}{k} \binom{42 - 6k}{14}
= \frac{10296957375}{6^{15}} = \frac{10296957375}{470184984576} \approx 0.0219 .
\]

\textbf{轉成計數}。 $15$ 次擲骰的結果共 $6^{15}$ 種，每種機率相同（資結書第~A.5~節）。 所求
機率是「$x_1 + \cdots + x_{15} = 43$、$1 \le x_i \le 6$ 的整數解個數」除以 $6^{15}$。 令
$y_i = x_i - 1$，變成
\[
y_1 + y_2 + \cdots + y_{15} = 43 - 15 = 28, \qquad 0 \le y_i \le 5 .
\]

\textbf{生成函數}。 每次擲骰的生成函數是 $1 + x + \cdots + x^5 = \dfrac{1 - x^6}{1 - x}$，解的個數是
\[
[x^{28}] \left(\frac{1 - x^6}{1 - x}\right)^{15}
= [x^{28}]\, (1 - x^6)^{15} \sum_{j \ge 0} \binom{j + 14}{14} x^j .
\]
$(1 - x^6)^{15} = \sum_k (-1)^k \binom{15}{k} x^{6k}$，和右邊的 $x^{28 - 6k}$ 配成 $x^{28}$，需要
$28 - 6k \ge 0$，即 $k \le 4$：
\[
\sum_{k=0}^{4} (-1)^k \binom{15}{k} \binom{28 - 6k + 14}{14}
= \sum_{k=0}^{4} (-1)^k \binom{15}{k} \binom{42 - 6k}{14} .
\]
這也就是排容原理：$\binom{42}{14}$ 是不管上限的解數，第 $k$ 項處理「有 $k$ 個 $y_i \ge 6$」。

\textbf{逐項計算}：
\begin{center}
\begin{tabular}{crrr}
\toprule
$k$ & $\binom{15}{k}$ & $\binom{42 - 6k}{14}$ & 帶符號的乘積 \\
\midrule
0 & 1 & 52860229080 & $+52860229080$ \\
1 & 15 & 3796297200 & $-56944458000$ \\
2 & 105 & 145422675 & $+15269380875$ \\
3 & 455 & 1961256 & $-892371480$ \\
4 & 1365 & 3060 & $+4176900$ \\
\midrule
 & & 合計 & $10296957375$ \\
\bottomrule
\end{tabular}
\end{center}
$6^{15} = 470184984576$，機率約 $0.0219$（化成最簡分數是 $\frac{1144106375}{52242776064}$）。

考場不能用計算機，寫到排容的和式並說明每一項的意義，通常就是要的答案；後面的數字是給讀者核對的。
附帶一提，$15$ 顆骰子的點數和以 $52.5$ 為中心對稱（$x_i \mapsto 7 - x_i$），所以點數和 $43$ 的
機率等於點數和 $62$ 的機率。
\end{pastsoln}
